\chapter*{Introduction}
\addcontentsline{toc}{chapter}{Introduction}
\VerdierPhysicalPageAnchor{0018}
\VerdierSemanticAnchor{intro}

\VerdierSemanticAnchor{intro-opening}
Nous proposons dans ce travail un formalisme de l'hyperhomologie.

\VerdierSemanticAnchor{intro-1}
\noindent\textbf{1.}\quad
Soient $\mathcal A$ et $\mathcal B$ deux catégories abéliennes,
$F:\mathcal A\to\mathcal B$ un foncteur additif. Lorsque la catégorie
$\mathcal A$ possède suffisamment d'objets injectifs, la construction des
foncteurs dérivés droits de $F$ est maintenant classique [1] : tout objet $X$
de $\mathcal A$ possède une résolution injective
$X\xrightarrow{\varepsilon}I^{\bullet}(X)$ unique à ``homotopie près''.
Les objets de cohomologie du complexe $F(I^{\bullet}(X))$ sont donc définis à
isomorphisme canonique près par l'objet $X$ et dépendent fonctoriellement de
$X$. Les foncteurs obtenus sont les foncteurs dérivés droits de $F$.

\VerdierSemanticAnchor{intro-1-homotopy-quasi}
Ainsi, on associe fonctoriellement à tout objet $X$ de $\mathcal A$ un complexe
d'objets de $\mathcal B$ défini ``à homotopie près'' dont on prend la
cohomologie. Cependant le plus souvent, les constructions usuelles de
l'algèbre homologique ne fournissent pas des complexes définis à homotopie
près, mais des complexes définis à \emph{quasi-isomorphisme près} (un
morphisme de complexes est appelé un quasi-isomorphisme s'il induit un
isomorphisme sur les objets de cohomologie). Pour illustrer ce point, donnons
deux exemples.

\VerdierSemanticAnchor{intro-1-f-acyclic}
Notons tout d'abord que, pour étudier les foncteurs dérivés de $F$, il y a lieu
de considérer des résolutions
$X\xrightarrow{\varepsilon'}\mathcal R(X)$ plus générales que les résolutions
injectives : les résolutions par des objets acycliques pour le foncteur $F$
(en théorie des faisceaux, résolutions par des faisceaux flasques, mous [23] ;
en théorie des modules, résolutions par des modules plats $\ldots$). Ces
résolutions ne sont plus nécessairement uniques à homotopie près. Mais étant
données une résolution $F$-acyclique
$X\xrightarrow{\varepsilon'}\mathcal R(X)$ et une résolution injective
$X\xrightarrow{\varepsilon}I^{\bullet}(X)$, il existe un et un seul (à
homotopie près) morphisme de résolutions, \emph{i.e.} un morphisme de complexes
$\rho:\mathcal R(X)\to I^{\bullet}(X)$ tel que le dia\discretionary{-}{}{}%
